\ifdefined\CORRECTION
\documentclass[correction]{td}
\else
\documentclass{td}
\fi

\inputexercisepath[src/corrections/TP]{src/exercises}

\usepackage[T1]{fontenc}
\usepackage[hidelinks]{hyperref}
\usepackage{amsmath,amssymb, mathrsfs, stmaryrd}
\newcommand\eqdef{\stackrel{\text{def}}{=}}

\usepackage{algos}
\usepackage{statemachines}
\usepackage{trees}

\codeUE{XLG4IU020}
\intituleUE{Langages et Automates}
\cursus{L2 Informatique}

\author{Matthieu \textsc{Perrin}}

\logo{src/img/logoUN.png}
\institution{Nantes Université}

\typeTP[5]
\title{Projet --- Enregistrements}

\hypersetup{
  pdftitle  = {Langages et Automates - Projet - Enregistrements},
  pdfauthor = {Matthieu Perrin},
  pdfsubject  = {Projet de L2 en Langages et Automates},
  pdfkeywords = {compilateur, lexer, parser, AST, typage, enregistrements}
}

\begin{document}

\maketitle

\vspace{-5mm}

\begin{center}
  \begin{minipage}{.9\textwidth}
    \it
    \textbf{Remarque importante.} Ce projet n'est pas à rendre. En revanche, il constitue l'entraînement principal pour le TP noté de fin de semestre.
    L'épreuve sera un exercice d'extension similaire à réaliser en temps limité. En particulier :
    \begin{itemize}
    \item Le code de départ sera identique à celui de ce projet, à l'exception de tests spécifiques.
    \item Le TP noté est un travail strictement individuel.
    \end{itemize}
  \end{minipage}
\end{center}

\vspace{3mm}

Le but de ce projet est d'ajouter au langage algorithmique développé dans les TPs précédents la notion d'\emph{enregistrement}.
Un enregistrement est un type structuré dont les instances sont des tuples nommés et typés.

Un type enregistrement est déclaré à l'aide du mot-clé \lstinline|enregistrement|, suivi d'un identifiant,
puis d'une liste de champs entre les mots-clés \lstinline|début| et \lstinline|fin|.
Chaque champ est constitué d'un nom, d'un type, et se termine par un point-virgule. Par exemple, le programme suivant déclare un enregistrement
\lstinline{Point} comme un couple d'entiers. 

\begin{lstlisting}[gobble=2]
  enregistrement Point début
    x : entier; y : entier;
  fin
\end{lstlisting}

Une instance d'un enregistrement est créée grâce au nom du type, suivi d'une liste de valeurs initiales entre parenthèses.
L'initialisation est positionnelle : les valeurs doivent être fournies dans l'ordre de déclaration des champs et respecter leur type.
Par exemple, l'expression \lstinline{Point(1,2)} crée un point pour lequel \lstinline{x} vaut $1$ et \lstinline{y} vaut $2$.

Les champs d'un enregistrement sont accessibles en lecture et en écriture à l'aide de l'opérateur point.
Cet opérateur est associatif à gauche et plus prioritaire que tous les autres opérateurs du langage.
Ainsi, l'expression \lstinline{c.centre.x} est interprétée comme \lstinline{((c.centre).x)}.

Lorsqu'une instance d'un enregistrement est convertie en chaîne de caractères,
par exemple lors d'un appel à \lstinline{écrire} ou d'une concaténation,
le résultat doit être le nom de l'enregistrement suivi de la liste de ses
champs entre parenthèses, dans l'ordre de déclaration,
comme l'expression utilisée pour créer l'instance.

Par exemple, le programme suivant doit produire la sortie ``Point(6,8) est sur le cercle Cercle(Point(3,4), 5)'' :

\begin{minipage}{.3\textwidth}
  \begin{lstlisting}[gobble=2]
    enregistrement Point
    début
      x : entier;
      y : entier;
    fin
    
    enregistrement Cercle
    début
      centre : Point;
      rayon  : entier;
    fin
  \end{lstlisting}
\end{minipage}%
\begin{minipage}{.7\textwidth}
  \begin{lstlisting}[gobble=2]
    algorithme
    variables
      p : Point;
      c : Cercle;
      d : Point;
    début
      p <- Point(6, 8);
      c <- Cercle(Point(3, 4), 5);
    
      d <- Point(c.centre.x - p.x, c.centre.y - p.y);
      si d.x * d.x + d.y * d.y = c.rayon * c.rayon alors
        écrire(p, " est sur le cercle ", c);
      fin si
    fin
  \end{lstlisting}
\end{minipage}%

\vspace{-4mm}
\begin{exercice}[Travail attendu]

  \begin{question}
  \item Clonez le dépôt du TP : \lstinline[language=bash]{git clone https://github.com/LangagesEtAutomates/lea-tp6-project}

    La compilation et l'exécution de l'interpréteur sont similaires aux TPs précédents.

  \item Implémentez la structure d'enregistrements dans le code déjà existant.
    Vous soignerez particulièrement l'analyse statique des types lors de la création de
    nouvelles instances et l'accès aux champs d'une instance existante. 
    
  \item Quelques tests sont déjà fournis pour vous aider à démarrer.
    N'hésitez pas à les adapter et à en ajouter en fonction de votre avancement. 
    
  \end{question}

\end{exercice}

\end{document}
